% !PRINTED 608
\weberpage{608}

%<<CONT
\weq{4}{P = \prod\nolimits \frac{1}{\left[1-N(\mathfrak{p}_1)^{-s}\right]},}
then according to (2) and (3), for $s=1$
\weq{5}{(s-1)\,P^n \text{ finite and different from zero,}}
and according to § 163, 6.:
\weq{6}{(s-1)\,P^h \text{ finite,}}
and if one eliminates $P$ from these, it follows that:
\[
(s-1)^{n-h} \text{ finite.}
\]
Therefore
\weq{7}{n \overline{\overline{>}} h.}
From this follows the theorem:

8. \gesperrt{The relative degree $n$ of the class field cannot be smaller than the class number $h$ of the group. If $n=h$, then $(s-1)P^h$ is not zero for $s=1$, and the sums $Q_2, Q_3, \ldots, Q_h$ are likewise different from zero.}

If an equation of degree $h$ can be formed in $\Omega$ from which a field can be derived that exhibits the characteristic properties of the class field, and if the degree of this field is not greater than $h$, then it follows that it cannot be smaller either, and so must be $= h$. With this the irreducibility of the equation in question is established.

On this rests the demonstration of the \gesperrt{irreducibility of the cyclotomic equation} which we gave in § 198 of the II. Band. There the field of rational numbers took the place of $\Omega$. In the next section we shall make further applications of this.

Whether there are groups for which the degree of the class field is greater than the class number remains undecided. I consider it improbable.

We prove further that there can be at most \gesperrt{one} class field. Every such field is generated by the adjunction of an algebraic number $\xi$ to $\Omega$.

Let $\xi$ be an \gesperrt{integer} generating the class field $\mathfrak{K} = \Omega(\xi)$, one which therefore satisfies an equation irreducible in $\Omega$
\weq{8}{f(\xi) = \xi^n + \alpha_1 \xi^{n-1} + \alpha_2 \xi^{n-2} + \cdots + \alpha_n = 0}
in which the $\alpha_1, \alpha_2, \ldots, \alpha_n$ are integers in $\Omega$.
